% Talnafræði — Þrenndir samliggjandi (teljandi lausn, án KLS/CRT)

% Lausn
Skrifum $2025=3^4\cdot5^2=81\cdot25$.
Í hverri þrennd $k(k+1)(k+2)$ er nákvæmlega ein tala margfeldi af $3$,
þannig að til að fá veldið $3^4$ þarf einhver talnanna að vera margfeldi af $81$.
Sömuleiðis þarf einhver talnanna að vera margfeldi af $25$.
Þannig að leitin snýst um þrenndir sem innihalda bæði
margfeldi af $81$ og margfeldi af $25$.

Látum $m$ vera margfeldi af $81$ á bilinu $[1,2027]$.
Þá er $m=81t$ með $1\le t\le 25$ (því $81\cdot 26>2027$).
Þrenndirnar sem innihalda $m$ eru
\[
\{m-2,m-1,m\},\quad \{m-1,m,m+1\},\quad \{m,m+1,m+2\}.
\]
Slík þrennd er gild þá og því aðeins að \emph{einhver} talnanna
$m-2,m-1,m,m+1,m+2$ sé margfeldi af $25$, sem jafngildir því að
\[
m\equiv 0,\ \pm 1,\ \pm 2 \pmod{25}. \tag{$\star$}
\]

Athugum að $81\equiv 6 \pmod{25}$, þannig að leifarnar
\(\mod 25\) fyrir $m=81t$ eru einfaldlega leifarnar fyrir
röðina $6,12,18,24, \ldots$ (þ.e.\ bætum 6 í hvert skref).
Í fyrsta $25$-töfluhringnum ($t=1,\ldots,25$) fáum við:
\[
\begin{array}{c|cccccccccc}
t       & 1&2&3&4&5&6&7&8&9&\cdots\\
81t\bmod 25 & 6&12&18&\mathbf{24}&5&11&17&\mathbf{23}&4&\cdots
\end{array}
\]
Haldið áfram til $t=25$ og athugið hvar leifarnar verða $0,\pm1,\pm2$:
\[
\boxed{t=4,\ 8,\ 17,\ 21,\ 25}
\]
sem samsvarar
\[
m=81t\in\{324,\ 648,\ 1377,\ 1701,\ 2025\}.
\]
Þetta eru einmitt þau $81$-margfeldi í bilinu $[1,2027]$ sem eru
í fjarlægð $\le 2$ frá margfeldum af $25$.

\medskip
Nú teljum við fyrir hvert $m$ hve mörg upphafsgildi $k$ gefa þrennd sem inniheldur
bæði $m$ fyrir $3^4$ og viðeigandi $25$-margfeldi:

\[
\begin{array}{c|c|l}
m \bmod 25 & 25\text{-margfeldið í nágrenninu} & \text{gild }k \text{ (fjöldi)}\\
\hline
0          & m                                  & k=m-2,\ m-1,\ m \quad (3)\\
1          & m-1                                & k=m-2,\ m-1 \quad (2)\\
-1\ (24)   & m+1                                & k=m-1,\ m \quad (2)\\
2          & m-2                                & k=m-2 \quad (1)\\
-2\ (23)   & m+2                                & k=m \quad (1)
\end{array}
\]

Fyrir okkar fimm $m$ gild eru leifarnar $24,\ 23,\ 2,\ 1,\ 0$,
svo heildarfjöldinn er
\[
2\;+\;1\;+\;1\;+\;2\;+\;3 \;=\; 9.
\]

Þar með er fjöldi slíkra þrennda á bilinu $k=1,\ldots,2025$ jafnt $9$.

Dæmið er einnig hægt að leysa með kínversku leifarsetningunni.